\documentclass[10pt,a4paper]{amsart}
\usepackage{amsmath,amssymb,amsthm,graphicx}
% Retained preamble for the Stacks Project chapter {CH} (shared project preamble).
% The chapter file pinned-original-source.tex only carries \input{{preamble}} on its first line,
% so no preamble text lives inside the chapter itself.  The Stacks Project uses one single shared
% preamble for every chapter: the snapshot stages it three times (sources/stacks-chow/preamble.tex,
% sources/stacks-etale-cohomology/preamble.tex and focused-seeds/AGP/sources/stacks-snapshot-20261011/
% preamble.tex) and those copies are identical in size.  The block below is copied mechanically,
% line for line, from sources/stacks-chow/preamble.tex lines 14-18 and 154-217: the Xy-pic setup and
% the native theorem declarations together with the mathematical macros.  Nothing was added, removed
% or redefined; the wrapper class and the amsmath/amssymb/amsthm/graphicx base packages are supplied
% by the assembler.
% For commutative diagrams we use Xy-pic
% Chapter: stacks-cohomology.tex -- Cohomology of Sheaves.
\usepackage[all]{xy}

% We use 2cell for 2-commutative diagrams.
\xyoption{2cell}

% Theorem environments.
%
\theoremstyle{plain}
\newtheorem{theorem}[subsection]{Theorem}
\newtheorem{proposition}[subsection]{Proposition}
\newtheorem{lemma}[subsection]{Lemma}

\theoremstyle{definition}
\newtheorem{definition}[subsection]{Definition}
\newtheorem{example}[subsection]{Example}
\newtheorem{exercise}[subsection]{Exercise}
\newtheorem{situation}[subsection]{Situation}

\theoremstyle{remark}
\newtheorem{remark}[subsection]{Remark}
\newtheorem{remarks}[subsection]{Remarks}

\numberwithin{equation}{subsection}

% Macros
%
\def\lim{\mathop{\mathrm{lim}}\nolimits}
\def\colim{\mathop{\mathrm{colim}}\nolimits}
\def\Spec{\mathop{\mathrm{Spec}}}
\def\Hom{\mathop{\mathrm{Hom}}\nolimits}
\def\Ext{\mathop{\mathrm{Ext}}\nolimits}
\def\SheafHom{\mathop{\mathcal{H}\!\mathit{om}}\nolimits}
\def\SheafExt{\mathop{\mathcal{E}\!\mathit{xt}}\nolimits}
\def\Sch{\mathit{Sch}}
\def\Mor{\mathop{\mathrm{Mor}}\nolimits}
\def\Ob{\mathop{\mathrm{Ob}}\nolimits}
\def\Sh{\mathop{\mathit{Sh}}\nolimits}
\def\NL{\mathop{N\!L}\nolimits}
\def\CH{\mathop{\mathrm{CH}}\nolimits}
\def\proetale{{pro\text{-}\acute{e}tale}}
\def\etale{{\acute{e}tale}}
\def\QCoh{\mathit{QCoh}}
\def\Ker{\mathop{\mathrm{Ker}}}
\def\Im{\mathop{\mathrm{Im}}}
\def\Coker{\mathop{\mathrm{Coker}}}
\def\Coim{\mathop{\mathrm{Coim}}}

% Boxtimes
%
\DeclareMathSymbol{\boxtimes}{\mathbin}{AMSa}{"02}

%
% Macros for moduli stacks/spaces
%
\def\QCohstack{\mathcal{QC}\!\mathit{oh}}
\def\Cohstack{\mathcal{C}\!\mathit{oh}}
\def\Spacesstack{\mathcal{S}\!\mathit{paces}}
\def\Quotfunctor{\mathrm{Quot}}
\def\Hilbfunctor{\mathrm{Hilb}}
\def\Curvesstack{\mathcal{C}\!\mathit{urves}}
\def\Polarizedstack{\mathcal{P}\!\mathit{olarized}}
\def\Complexesstack{\mathcal{C}\!\mathit{omplexes}}
% \Pic is the operator that assigns to X its picard group, usage \Pic(X)
% \Picardstack_{X/B} denotes the Picard stack of X over B
% \Picardfunctor_{X/B} denotes the Picard functor of X over B
\def\Pic{\mathop{\mathrm{Pic}}\nolimits}
\def\Picardstack{\mathcal{P}\!\mathit{ic}}
\def\Picardfunctor{\mathrm{Pic}}
\def\Deformationcategory{\mathcal{D}\!\mathit{ef}}
\usepackage{cleveref}
\crefname{originallemma}{original lemma}{original lemmas}
\Crefname{originallemma}{Original lemma}{Original lemmas}
\makeatletter
\newlabel{homology-lemma-ss-converges-trivial}{{0BK5}{source}{}{sourcepointer-1}{}}
\newlabel{homology-lemma-ss-converges-trivial@cref}{{[originallemma][0][]0BK5}{[1][0][]source}{}{}{}}
\makeatother
\crefname{originalsection}{original section}{original sections}
\Crefname{originalsection}{Original section}{Original sections}
\makeatletter
\newlabel{homology-section-filtered-complex}{{012K}{source}{}{sourcepointer-2}{}}
\newlabel{homology-section-filtered-complex@cref}{{[originalsection][0][]012K}{[1][0][]source}{}{}{}}
\makeatother
\crefname{originallemma}{original lemma}{original lemmas}
\Crefname{originallemma}{Original lemma}{Original lemmas}
\makeatletter
\newlabel{injectives-lemma-K-injective-embedding-filtration}{{0BKI}{source}{}{sourcepointer-3}{}}
\newlabel{injectives-lemma-K-injective-embedding-filtration@cref}{{[originallemma][0][]0BKI}{[1][0][]source}{}{}{}}
\makeatother
\title{Spectral sequence of a filtered complex of modules on a ringed space}
\author{The Stacks Project Authors}
\date{stacks-project:cohomology.tex:local-snapshot-20261011; GNU FDL 1.2 or later}
\begin{document}
\maketitle
\noindent\textit{Editorial note.} All delivered mathematical text is the author's own, copied from the pinned Stacks Project chapter file cohomology.tex; no statement or proof step was rewritten, reordered or removed.\par\medskip
\noindent\textit{Editorial note.} The proof opens with the author's own parenthetical pointing to a remark for the bounded below case with finite filtrations; the parenthetical and that reference are reproduced as written and carried as a navigation pointer.\par\medskip
\noindent\textit{Editorial note.} The counted spans carry 3 navigation stubs and no bibliography citation.\par\medskip
\noindent\textit{Editorial note.} The preamble is the chapter's retained minimal preamble (sources/stacks-cohomology/retained-preamble.tex lines 1-81), a verbatim excerpt of the native preamble.\par\medskip
\noindent\textit{Editorial note.} The label lemma-spectral-sequence-filtered-object is defined only in cohomology.tex, so no cross-chapter restatement had to be weighed for this card.\par\medskip
\noindent\textit{Allowed standard prerequisites (editorial).}\begin{itemize}
\item Embedding of a filtered complex into a filtered K-injective complex. Injectives, Lemma injectives-lemma-K-injective-embedding-filtration, invoked at proof lines 6843-6846. Supplies the representative on which the spectral sequence is constructed.
\item Spectral sequence of a filtered complex. Homology, Section homology-section-filtered-complex, invoked at proof lines 6847-6854. Supplies the construction of the spectral sequence from the filtered complex of global sections.
\item Boundedness and convergence of the spectral sequence under the trivial vanishing conditions. Homology, Lemma homology-lemma-ss-converges-trivial, invoked at proof lines 6857-6858. Supplies the convergence statement of the second half of the lemma.
\end{itemize}\medskip
\ifdefined\hypertarget\hypertarget{sourcepointer-1}{}\fi
\noindent\textit{Source reference.} Original lemma 0BK5: sources/stacks-chow/tags.txt, source lines 14548--14548. Pointer outside the counted span, invoked once for context; the load-bearing argument of the counted statement is the author's proof reproduced in full in the proof span.\par\medskip
\ifdefined\hypertarget\hypertarget{sourcepointer-2}{}\fi
\noindent\textit{Source reference.} Original section 012K: sources/stacks-chow/tags.txt, source lines 1402--1402. Pointer outside the counted span, invoked once for context; the load-bearing argument of the counted statement is the author's proof reproduced in full in the proof span.\par\medskip
\ifdefined\hypertarget\hypertarget{sourcepointer-3}{}\fi
\noindent\textit{Source reference.} Original lemma 0BKI: sources/stacks-chow/tags.txt, source lines 14561--14561. Pointer outside the counted span, invoked once for context; the load-bearing argument of the counted statement is the author's proof reproduced in full in the proof span.\par\medskip
\begin{lemma}
\label{lemma-spectral-sequence-filtered-object}
Let $(X, \mathcal{O}_X)$ be a ringed space. Let $\mathcal{F}^\bullet$ be a
filtered complex of $\mathcal{O}_X$-modules. There exists a canonical
spectral sequence $(E_r, \text{d}_r)_{r \geq 1}$ of bigraded
$\Gamma(X, \mathcal{O}_X)$-modules with $d_r$ of bidegree $(r, -r + 1)$ and
$$
E_1^{p, q} = H^{p + q}(X, \text{gr}^p\mathcal{F}^\bullet)
$$
If for every $n$ we have
$$
H^n(X, F^p\mathcal{F}^\bullet) = 0\text{ for }p \gg 0
\quad\text{and}\quad
H^n(X, F^p\mathcal{F}^\bullet) = H^n(X, \mathcal{F}^\bullet)\text{ for }p \ll 0
$$
then the spectral sequence is bounded and converges to
$H^*(X, \mathcal{F}^\bullet)$.
\end{lemma}

\begin{proof}
(For a proof in case the complex is a bounded below complex
of modules with finite filtrations, see the remark below.)
Choose an map of filtered complexes
$j : \mathcal{F}^\bullet \to \mathcal{J}^\bullet$ as in
Injectives, Lemma
\ref{injectives-lemma-K-injective-embedding-filtration}.
The spectral sequence is the spectral sequence of
Homology, Section \ref{homology-section-filtered-complex}
associated to the filtered complex
$$
\Gamma(X, \mathcal{J}^\bullet)
\quad\text{with}\quad
F^p\Gamma(X, \mathcal{J}^\bullet) = \Gamma(X, F^p\mathcal{J}^\bullet)
$$
Since cohomology is computed by evaluating on K-injective representatives
we see that the $E_1$ page is as stated in the lemma.
The convergence and boundedness under the stated conditions
follows from Homology, Lemma \ref{homology-lemma-ss-converges-trivial}.
\end{proof}

\end{document}
